\documentclass[10pt,a4paper]{amsart}
\usepackage{amsmath,amssymb,amsthm,graphicx}
\usepackage{verbatim}
\newenvironment{reference}{\comment}{\endcomment}
%\newenvironment{reference}{}{}
\newenvironment{slogan}{\comment}{\endcomment}
\newenvironment{history}{\comment}{\endcomment}

% For commutative diagrams we use Xy-pic
\usepackage[all]{xy}

% We use 2cell for 2-commutative diagrams.
\xyoption{2cell}
% Theorem environments.
%
\theoremstyle{plain}
\newtheorem{theorem}[subsection]{Theorem}
\newtheorem{proposition}[subsection]{Proposition}
\newtheorem{lemma}[subsection]{Lemma}

\theoremstyle{definition}
\newtheorem{definition}[subsection]{Definition}
\newtheorem{example}[subsection]{Example}
\newtheorem{exercise}[subsection]{Exercise}
\newtheorem{situation}[subsection]{Situation}

\theoremstyle{remark}
\newtheorem{remark}[subsection]{Remark}
\newtheorem{remarks}[subsection]{Remarks}

\numberwithin{equation}{subsection}

% Macros
%
\def\lim{\mathop{\mathrm{lim}}\nolimits}
\def\colim{\mathop{\mathrm{colim}}\nolimits}
\def\Spec{\mathop{\mathrm{Spec}}}
\def\Hom{\mathop{\mathrm{Hom}}\nolimits}
\def\Ext{\mathop{\mathrm{Ext}}\nolimits}
\def\SheafHom{\mathop{\mathcal{H}\!\mathit{om}}\nolimits}
\def\SheafExt{\mathop{\mathcal{E}\!\mathit{xt}}\nolimits}
\def\Sch{\mathit{Sch}}
\def\Mor{\mathop{\mathrm{Mor}}\nolimits}
\def\Ob{\mathop{\mathrm{Ob}}\nolimits}
\def\Sh{\mathop{\mathit{Sh}}\nolimits}
\def\NL{\mathop{N\!L}\nolimits}
\def\CH{\mathop{\mathrm{CH}}\nolimits}
\def\proetale{{pro\text{-}\acute{e}tale}}
\def\etale{{\acute{e}tale}}
\def\QCoh{\mathit{QCoh}}
\def\Ker{\mathop{\mathrm{Ker}}}
\def\Im{\mathop{\mathrm{Im}}}
\def\Coker{\mathop{\mathrm{Coker}}}
\def\Coim{\mathop{\mathrm{Coim}}}

% Boxtimes
%
\DeclareMathSymbol{\boxtimes}{\mathbin}{AMSa}{"02}

%
% Macros for moduli stacks/spaces
%
\def\QCohstack{\mathcal{QC}\!\mathit{oh}}
\def\Cohstack{\mathcal{C}\!\mathit{oh}}
\def\Spacesstack{\mathcal{S}\!\mathit{paces}}
\def\Quotfunctor{\mathrm{Quot}}
\def\Hilbfunctor{\mathrm{Hilb}}
\def\Curvesstack{\mathcal{C}\!\mathit{urves}}
\def\Polarizedstack{\mathcal{P}\!\mathit{olarized}}
\def\Complexesstack{\mathcal{C}\!\mathit{omplexes}}
% \Pic is the operator that assigns to X its picard group, usage \Pic(X)
% \Picardstack_{X/B} denotes the Picard stack of X over B
% \Picardfunctor_{X/B} denotes the Picard functor of X over B
\def\Pic{\mathop{\mathrm{Pic}}\nolimits}
\def\Picardstack{\mathcal{P}\!\mathit{ic}}
\def\Picardfunctor{\mathrm{Pic}}
\def\Deformationcategory{\mathcal{D}\!\mathit{ef}}
\usepackage{cleveref}
\crefname{originalreference}{original reference}{original references}
\Crefname{originalreference}{Original reference}{Original references}
\makeatletter
\newlabel{lemma-compute-what-it-should-be}{{0CTD}{source}{}{sourcepointer-1}{}}
\newlabel{lemma-compute-what-it-should-be@cref}{{[originalreference][0][]0CTD}{[1][0][]source}{}{}{}}
\makeatother
\crefname{originallemma}{original lemma}{original lemmas}
\Crefname{originallemma}{Original lemma}{Original lemmas}
\makeatletter
\newlabel{more-algebra-lemma-Rlim-perfect-gives-perfect}{{0CQG}{source}{}{sourcepointer-2}{}}
\newlabel{more-algebra-lemma-Rlim-perfect-gives-perfect@cref}{{[originallemma][0][]0CQG}{[1][0][]source}{}{}{}}
\makeatother
\crefname{originallemma}{original lemma}{original lemmas}
\Crefname{originallemma}{Original lemma}{Original lemmas}
\makeatletter
\newlabel{sites-cohomology-lemma-dual-perfect-complex}{{08JJ}{source}{}{sourcepointer-3}{}}
\newlabel{sites-cohomology-lemma-dual-perfect-complex@cref}{{[originallemma][0][]08JJ}{[1][0][]source}{}{}{}}
\makeatother
\crefname{originalsection}{original section}{original sections}
\Crefname{originalsection}{Original section}{Original sections}
\makeatletter
\newlabel{sites-cohomology-section-global-RHom}{{0B6E}{source}{}{sourcepointer-4}{}}
\newlabel{sites-cohomology-section-global-RHom@cref}{{[originalsection][0][]0B6E}{[1][0][]source}{}{}{}}
\makeatother
\crefname{originalreference}{original reference}{original references}
\Crefname{originalreference}{Original reference}{Original references}
\makeatletter
\newlabel{situation-existence}{{08B7}{source}{}{sourcepointer-5}{}}
\newlabel{situation-existence@cref}{{[originalreference][0][]08B7}{[1][0][]source}{}{}{}}
\makeatother
\crefname{originallemma}{original lemma}{original lemmas}
\Crefname{originallemma}{Original lemma}{Original lemmas}
\makeatletter
\newlabel{spaces-perfect-lemma-base-change-RHom-perfect}{{0A1R}{source}{}{sourcepointer-6}{}}
\newlabel{spaces-perfect-lemma-base-change-RHom-perfect@cref}{{[originallemma][0][]0A1R}{[1][0][]source}{}{}{}}
\makeatother
\title{Computing the derived global sections and derived hom against a perfect object along the descent tower}
\author{The Stacks Project Authors}
\date{stacks-project:spaces-flat.tex:local-snapshot-20261012; GNU FDL 1.2 or later}
\begin{document}
\maketitle
\noindent\textit{Editorial note.} All delivered mathematical text is the author's own, copied from the pinned Stacks Project chapter file spaces-flat.tex; no statement or proof step was rewritten, reordered or removed.\par\medskip
\noindent\textit{Editorial note.} The counted proof is the author's complete argument: a stage-wise RHom computation with base change compatible perfect transitions, the derived-limit commutation of RHom\_X(E, -), and a duality translation from the hom statement to the tensor statement\par\medskip
\noindent\textit{Editorial note.} The label lemma-compute-against-perfect names the counted statement, and the proof environment immediately following it in the source is delivered verbatim as its proof.\par\medskip
\noindent\textit{Editorial note.} The counted statement is carried by the author's own lemma environment in the chapter source.\par\medskip
\noindent\textit{Allowed standard prerequisites (editorial).}\begin{itemize}
\item The descent tower situation with its system of objects and the object K defined as the derived limit of the system, and the identification of RHom with RGamma composed with the internal Hom sheaf
\item The base change lemma for RHom against a perfect object: the stages RHom\_\{X\_n\}(E\_n, F\_n) are perfect over A\_n with formation commuting with base change
\item The More on Algebra lemma that a derived limit of perfect objects with compatible base change data is perfect and recovers the stages on base change
\item The duality between derived global sections and derived hom translating tensor statements into hom statements for perfect objects
\end{itemize}\medskip
\ifdefined\hypertarget\hypertarget{sourcepointer-1}{}\fi
\noindent\textit{Source reference.} Original reference 0CTD: sources/stacks-spaces-flat/pinned-original-source.tex, source lines 2388--2388. Prerequisite pointer invoked at proof line(s) 2407; it names a standard fact of the same chapter that the counted argument uses, and every step of the counted proof built on it is retained verbatim in the delivered body.\par\medskip
\ifdefined\hypertarget\hypertarget{sourcepointer-2}{}\fi
\noindent\textit{Source reference.} Original lemma 0CQG: sources/stacks-chow/tags.txt, source lines 15961--15961. Prerequisite pointer invoked at proof line(s) 2436; it names a standard fact of a sibling chapter that the counted argument uses, and every step of the counted proof built on it is retained verbatim in the delivered body.\par\medskip
\ifdefined\hypertarget\hypertarget{sourcepointer-3}{}\fi
\noindent\textit{Source reference.} Original lemma 08JJ: sources/stacks-chow/tags.txt, source lines 10852--10852. Prerequisite pointer invoked at proof line(s) 2451; it names a standard fact of a sibling chapter that the counted argument uses, and every step of the counted proof built on it is retained verbatim in the delivered body.\par\medskip
\ifdefined\hypertarget\hypertarget{sourcepointer-4}{}\fi
\noindent\textit{Source reference.} Original section 0B6E: sources/stacks-chow/tags.txt, source lines 14067--14067. Prerequisite pointer invoked at proof line(s) 2429; it names a standard fact of a sibling chapter that the counted argument uses, and every step of the counted proof built on it is retained verbatim in the delivered body.\par\medskip
\ifdefined\hypertarget\hypertarget{sourcepointer-5}{}\fi
\noindent\textit{Source reference.} Original reference 08B7: sources/stacks-spaces-flat/pinned-original-source.tex, source lines 2362--2362. Prerequisite pointer invoked at proof line(s) 2406; it names a standard fact of the same chapter that the counted argument uses, and every step of the counted proof built on it is retained verbatim in the delivered body.\par\medskip
\ifdefined\hypertarget\hypertarget{sourcepointer-6}{}\fi
\noindent\textit{Source reference.} Original lemma 0A1R: sources/stacks-chow/tags.txt, source lines 12679--12679. Prerequisite pointer invoked at proof line(s) 2431; it names a standard fact of a sibling chapter that the counted argument uses, and every step of the counted proof built on it is retained verbatim in the delivered body.\par\medskip
\begin{lemma}
\label{lemma-compute-against-perfect}
In Situation \ref{situation-existence} let $K$ be as in
Lemma \ref{lemma-compute-what-it-should-be}. For any perfect
object $E$ of $D(\mathcal{O}_X)$ we have
\begin{enumerate}
\item $M = R\Gamma(X, K \otimes^\mathbf{L} E)$ is a perfect object of $D(A)$
and there is a canonical isomorphism
$R\Gamma(X_n, \mathcal{F}_n \otimes^\mathbf{L} E|_{X_n}) =
M \otimes_A^\mathbf{L} A_n$
in $D(A_n)$,
\item $N = R\Hom_X(E, K)$ is a perfect object of $D(A)$
and there is a canonical isomorphism
$R\Hom_{X_n}(E|_{X_n}, \mathcal{F}_n) = N \otimes_A^\mathbf{L} A_n$
in $D(A_n)$.
\end{enumerate}
In both statements $E|_{X_n}$ denotes the derived pullback
of $E$ to $X_n$.
\end{lemma}

\begin{proof}
Proof of (2). Write $E_n = E|_{X_n}$ and
$N_n = R\Hom_{X_n}(E_n, \mathcal{F}_n)$.
Recall that $R\Hom_{X_n}(-, -)$ is equal to
$R\Gamma(X_n, R\SheafHom(-, -))$, see
Cohomology on Sites, Section \ref{sites-cohomology-section-global-RHom}.
Hence by Derived Categories of Spaces, Lemma
\ref{spaces-perfect-lemma-base-change-RHom-perfect}
we see that $N_n$ is a perfect object of $D(A_n)$
whose formation commutes with base change. Thus the maps
$N_n \otimes_{A_n}^\mathbf{L} A_{n - 1} \to N_{n - 1}$
coming from $\varphi_n$ are isomorphisms.
By More on Algebra, Lemma \ref{more-algebra-lemma-Rlim-perfect-gives-perfect}
we find that $R\lim N_n$ is perfect and
that its base change back to $A_n$ recovers $N_n$.
On the other hand, the exact functor
$R\Hom_X(E, -) : D_\QCoh(\mathcal{O}_X) \to D(A)$
of triangulated categories commutes with products
and hence with derived limits, whence
$$
R\Hom_X(E, K) =
R\lim R\Hom_X(E, \mathcal{F}_n) =
R\lim R\Hom_X(E_n, \mathcal{F}_n) =
R\lim N_n
$$
This proves (2). To see that (1) holds, translate it into (2)
using Cohomology on Sites, Lemma
\ref{sites-cohomology-lemma-dual-perfect-complex}.
\end{proof}

\end{document}
